% Job posting: https://ca.linkedin.com/jobs/view/scientifique-en-traitement-de-signal-signal-processing-scientist-at-flyscan-systems-inc-4434616270
\documentclass[11pt]{article}
\usepackage[left=1in,top=1in,right=1in,bottom=1in]{geometry}
\usepackage[hidelinks]{hyperref}
\usepackage{setspace}
\usepackage[T1]{fontenc}
\usepackage{newtxtext,newtxmath}
\ifdefined\pdfgentounicode
  \input{glyphtounicode}
  \pdfgentounicode=1
\fi
\setstretch{1.1}
\pagenumbering{gobble}
\begin{document}
\pagestyle{empty}

\begin{flushleft}
Nishal Stanislaus Silva, PhD \\
Toronto, ON, Canada \\
nishal.silva@hotmail.com \,|\, +1 613 200 2030 \,|\, \href{https://nishal.xyz}{nishal.xyz} \,|\, \href{https://www.linkedin.com/in/nishal-silva}{linkedin.com/in/nishal-silva}
\end{flushleft}

\vspace{1em}

Dear Hiring Manager,

\vspace{0.5em}

My MSc thesis at Sheffield Hallam University applied the Stockwell (S-)transform -- a time-frequency signal-processing technique -- to the problem of extracting reliable transient onset events from noisy, continuous audio recordings, published at IEEE Asilomar in 2018. That problem, at its core, is the same one Flyscan's condition-monitoring instrumentation solves every day: separating a real, physically meaningful signal from a noisy continuous sensor stream, reliably and with a defensible margin of confidence. It is the reason your Signal Processing Scientist posting caught my attention immediately, and the reason I believe my training maps directly onto it rather than adjacently.

That early work became the foundation for the signal-processing practice I have built since. As a postdoctoral researcher at the University of Trento, I owned a full DSP feature-engineering pipeline end to end -- from raw streaming sensor acquisition through STFT/onset-detection feature design to a deployed real-time detector running at 14 ms inference on embedded ARM hardware, 74x faster than the classical baseline it replaced. What I would highlight for this role specifically is the benchmark discipline behind that result: I designed frozen-protocol test suites and ablation studies that quantified exactly which features carried signal versus noise, and what each accuracy/latency/CPU trade-off cost, including a written record of what did not work. Earlier, at MAS Holdings, I spent five and a half years shipping C++ and Python signal- and image-processing systems onto embedded and PLC-integrated production hardware at manufacturing scale -- the same real-time deployment discipline that turns a working algorithm into a field-reliable instrument.

One honest note on language: I am building toward Quebec City as a fully engaged candidate, and my French is currently at a basic, actively-improving level (approximately CEFR A2) -- not yet a working professional level. I would welcome clarity on whether the role requires daily working French for the day-to-day technical work, or whether English is sufficient there while I continue building my French; either way, I am committed to keeping up that progress and to relocating to Quebec City for the right opportunity.

I am a Canadian Permanent Resident, eligible to work anywhere in Canada without sponsorship, and available immediately. I would welcome the opportunity to discuss how my signal-processing background could contribute to Flyscan's condition-monitoring platform.

\vspace{1em}

Sincerely, \\
Nishal Stanislaus Silva, PhD

\end{document}
